\documentclass[tikz,border=5mm]{standalone}
\usepackage{tikz}
\usetikzlibrary{shapes.geometric,shapes.misc,arrows.meta,positioning,
                fit,backgrounds,calc,decorations.pathreplacing,decorations.markings}
\usepackage[utf8]{vietnam}

% ── Bảng màu (mỗi giai đoạn một màu) ──
\definecolor{cBlue}{RGB}{66,133,244}
\definecolor{cBlueFill}{RGB}{225,237,255}
\definecolor{cTeal}{RGB}{0,151,167}
\definecolor{cTealFill}{RGB}{220,246,249}

\definecolor{p1}{RGB}{76,175,80}     % Khởi tạo  — xanh lá
\definecolor{p1F}{RGB}{225,245,228}
\definecolor{p2}{RGB}{230,162,60}    % Kết nối   — cam
\definecolor{p2F}{RGB}{255,243,224}
\definecolor{p3}{RGB}{126,87,194}    % Mô phỏng  — tím
\definecolor{p3F}{RGB}{237,231,246}
\definecolor{p4}{RGB}{211,84,85}     % Phân tích — đỏ
\definecolor{p4F}{RGB}{255,232,234}
\definecolor{p5}{RGB}{0,151,167}     % Tối ưu    — xanh ngọc
\definecolor{p5F}{RGB}{220,246,249}

\tikzset{
    % Trong circle CHỈ chứa số (không có chữ) ⇒ không bao giờ tràn.
    phasenode/.style={circle, draw=#1, fill=white, very thick,
                      minimum size=1.4cm, inner sep=1pt,
                      font=\Large\bfseries, align=center},
    % Hộp mô tả bên ngoài circle: chứa cả tên giai đoạn lẫn diễn giải.
    phaselabel/.style={rectangle, rounded corners=4pt, draw=#1,
                       fill=white, thick,
                       minimum width=3.6cm, minimum height=1.2cm,
                       font=\footnotesize, align=center, inner sep=4pt},
    centerbox/.style={rectangle, rounded corners=5pt, very thick,
                      minimum width=3.0cm, minimum height=1.10cm,
                      font=\small\bfseries, align=center,
                      inner sep=4pt},
    cycarr/.style={-{Stealth[length=3.5mm,width=2.6mm]}, very thick, #1},
    bidir/.style={{Stealth[length=2.8mm,width=2.2mm]}-{Stealth[length=2.8mm,width=2.2mm]},
                  very thick, #1},
}

\begin{document}
\begin{tikzpicture}

% ====================================================================
%  TRUNG TÂM — Thực thể vật lý  <-->  Mô hình số
%  Tách hai hộp ra xa và đặt nhãn nối ở khoảng giữa thật rộng
% ====================================================================
\node[centerbox, draw=cBlue, fill=cBlueFill] (pe) at (-3.8, 0)
    {Physical Entity (PE)\\[-1pt]{\scriptsize\mdseries Mạng / thiết bị IoT}};

\node[centerbox, draw=cTeal, fill=cTealFill] (dm) at ( 3.8, 0)
    {Digital Model (DM)\\[-1pt]{\scriptsize\mdseries Bản sao ảo}};

% Mũi tên hai chiều "Data Connection" giữa hai hộp
\draw[bidir=black!60] (pe.east) -- (dm.west);

% Nhãn "Data Connection" – đặt phía trên đoạn nối, có viền nhẹ để nổi bật
\node[font=\scriptsize\bfseries, fill=white,
      draw=black!30, semithick, rounded corners=2pt, inner sep=3pt]
      at (0, 0.45) {Data Connection (DC)};
% Phụ đề bên dưới đoạn nối
\node[font=\scriptsize\itshape, text=black!60, fill=white, inner sep=2pt]
      at (0, -0.30) {Telemetry $\,\leftrightarrows\,$ Control commands};

% ====================================================================
%  VÒNG NGOÀI — 5 giai đoạn xếp trên ngũ giác (thuận chiều kim đồng hồ)
% ====================================================================
\def\R{6.2}      % bán kính vòng giai đoạn (tăng để có chỗ cho hộp giữa)

% --- Giai đoạn 1: Create ---
\path (90:\R)   coordinate (cP1);
\node[phasenode=p1, fill=p1F] at (cP1) {1};
\node[phaselabel=p1, above=0.55cm of cP1]
      {\textbf{\color{p1}Create}\\Khởi tạo Digital Model\\từ đặc tả Physical Entity};

% --- Giai đoạn 2: Connect ---
\path (18:\R)   coordinate (cP2);
\node[phasenode=p2, fill=p2F] at (cP2) {2};
\node[phaselabel=p2, right=0.55cm of cP2]
      {\textbf{\color{p2}Connect}\\Thiết lập Data Connection\\trao đổi dữ liệu hai chiều};

% --- Giai đoạn 3: Simulate ---
\path (-54:\R)  coordinate (cP3);
\node[phasenode=p3, fill=p3F] at (cP3) {3};
\node[phaselabel=p3, below right=0.30cm and 0.55cm of cP3]
      {\textbf{\color{p3}Simulate}\\Mô phỏng hành vi của PE\\trong sandbox an toàn};

% --- Giai đoạn 4: Analyze ---
\path (-126:\R) coordinate (cP4);
\node[phasenode=p4, fill=p4F] at (cP4) {4};
\node[phaselabel=p4, below left=0.30cm and 0.55cm of cP4]
      {\textbf{\color{p4}Analyze}\\Phát hiện bất thường\\và đánh giá hành vi};

% --- Giai đoạn 5: Optimize \& Actuate ---
\path (162:\R)  coordinate (cP5);
\node[phasenode=p5, fill=p5F] at (cP5) {5};
\node[phaselabel=p5, left=0.55cm of cP5]
      {\textbf{\color{p5}Optimize \& Actuate}\\Phản hồi tri thức để\\điều khiển PE};

% ====================================================================
%  MŨI TÊN VÒNG (thuận chiều kim đồng hồ)
% ====================================================================
\def\gap{14}     % khoảng cách (độ) tránh node
\foreach \fromAng/\toAng/\colr in {%
    90/18/p1,
    18/-54/p2,
    -54/-126/p3,
    -126/-234/p4,
    -234/-270/p5}
{
    \draw[cycarr=\colr]
      (\fromAng-\gap:\R) arc (\fromAng-\gap:\toAng+\gap:\R);
}

% ====================================================================
%  ĐƯỜNG CHẤM — Trung tâm  <->  Vòng đời
% ====================================================================
\draw[->, thick, cBlue!70, dashed] (pe.north) to[bend left=15] (cP1);
\draw[->, thick, p5!80, dashed]    (cP5) to[bend left=15] (pe.west);
\draw[->, thick, p3!80, dashed]    (cP3) to[bend left=15] (dm.south east);
\draw[->, thick, p2!80, dashed]    (cP2) to[bend left=15] (dm.north east);

% ====================================================================
%  TIÊU ĐỀ + PHỤ ĐỀ
% ====================================================================
\node[font=\Large\bfseries, text=black!75]
      at (0, \R + 2.6) {Vòng đời Bản sao số (Digital Twin Lifecycle)};
\node[font=\footnotesize\itshape, text=black!55]
      at (0, \R + 2.0)
      {Vòng lặp liên tục: Create $\to$ Connect $\to$ Simulate $\to$ Analyze $\to$ Optimize \& Actuate};

% ====================================================================
%  CHÚ THÍCH (đáy) — dùng MŨI TÊN minh họa thay vì text "mũi tên liền"
% ====================================================================
\begin{scope}[shift={(0, -\R - 2.0)}]
    % Khung chú thích — đủ rộng để chứa cả 3 mục + nhãn cuối
    \node[draw=black!40, rounded corners=4pt, fill=black!4,
          minimum width=18cm, minimum height=1.2cm,
          inner xsep=10pt, inner ysep=8pt] (legbox) {};

    \node[font=\small\bfseries, anchor=west]
          at ($(legbox.west)+(0.3,0)$) (lblTitle) {Chú thích:};

    % --- Mục 1: mũi tên liền = chuyển giai đoạn ---
    \draw[cycarr=p1!80] ($(lblTitle.east)+(0.4,0)$) -- ++(0.8,0) coordinate (a1end);
    \node[font=\scriptsize, anchor=west, text=black!75] at ($(a1end)+(0.10,0)$)
          (txt1) {chuyển giai đoạn};

    % --- Mục 2: mũi tên đứt = liên kết thực thể ↔ giai đoạn ---
    \draw[->, thick, dashed, cBlue!80]
          ($(txt1.east)+(0.5,0)$) -- ++(0.8,0) coordinate (a2end);
    \node[font=\scriptsize, anchor=west, text=black!75] at ($(a2end)+(0.10,0)$)
          (txt2) {liên kết thực thể $\leftrightarrow$ giai đoạn};

    % --- Mục 3: mũi tên hai chiều = đồng bộ thời gian thực ---
    \draw[bidir=black!70]
          ($(txt2.east)+(0.5,0)$) -- ++(0.8,0) coordinate (a3end);
    \node[font=\scriptsize, anchor=west, text=black!75] at ($(a3end)+(0.10,0)$)
          {đồng bộ thời gian thực};
\end{scope}

% ====================================================================
%  NỀN HÌNH TRÒN (chỉ tô nhẹ, không vẽ viền dashed để khỏi xuyên qua node)
% ====================================================================
\begin{scope}[on background layer]
    \fill[black!2] (0,0) circle (\R);
\end{scope}

\end{tikzpicture}
\end{document}







